\addcontentsline{toc}{subsubsection}{Solution 62. The Hidden Square Root}

\begin{dlsummary}[title={\bfseries Solution 62. The Hidden Square Root},breakable]

\begin{proof}
The unique solution is
\[
\boxed{f(x)=\sqrt{|x|}\qquad(x\in\mathbb Q).}
\]

\paragraph{1. Turn the multiplicative inequality into equality.}

For \(x\ne0\), supermultiplicativity gives
\[
1=f(1)\ge f(x)f(1/x).
\]
This repeats the stated bound, so we need a stronger argument.

Apply supermultiplicativity to \(4x\) and \(1/x\):
\[
2=f(4)\ge f(4x)f(1/x)\ge2f(x)f(1/x).
\]
Again, this only supplies the same upper bound.

In fact, these hypotheses do \textbf{not} force multiplicativity: the reciprocal condition points in the same direction as supermultiplicativity at a reciprocal pair.

\textbf{The corrected problem replaces the reciprocal condition by}
\[
\boxed{\qquad f(x)f(1/x)\ge1\qquad(x\ne0).\qquad} \tag{3}
\]

With this correction, the two bounds force
\[
f(x)f(1/x)=1.
\]
In particular, \(f(x)>0\) for \(x\ne0\).

Now, for nonzero \(x,y\),
\[
f(x)=f((xy)(1/y))
\ge f(xy)f(1/y)=\frac{f(xy)}{f(y)}.
\]
Combined with supermultiplicativity, this proves
\[
\boxed{f(xy)=f(x)f(y).} \tag{4}
\]

Consequently,
\[
f(2)=\sqrt2,\qquad f(-1)=1,\qquad f(-x)=f(x).
\]

\paragraph{2. Prove an upper bound using binary expansions.}

Write a positive integer \(N\) in binary:
\[
N=\sum_{j=0}^{s}\varepsilon_j2^j,
\qquad \varepsilon_j\in\{0,1\},\quad 2^s\le N.
\]
Subadditivity and multiplicativity give
\[
f(N)\le\sum_{j=0}^{s}(\sqrt2)^j
< C\sqrt N,
\qquad
C=\frac{\sqrt2}{\sqrt2-1}.
\]
Apply this to \(N=n^r\):
\[
f(n)^r=f(n^r)<C\,n^{r/2}.
\]
Taking \(r\)-th roots and letting \(r\to\infty\),
\[
\boxed{f(n)\le\sqrt n.} \tag{5}
\]

\paragraph{3. Prove the reverse bound using base-\(b\) expansions.}

Fix an integer \(b\ge2\), and write \(A=f(b)>0\). Let
\[
M=\max\{f(0),f(1),\ldots,f(b-1)\}.
\]

Write \(2^k\) in base \(b\):
\[
2^k=d_0+d_1b+\cdots+d_sb^s,
\qquad 0\le d_j<b,\quad b^s\le2^k.
\]
Then
\[
2^{k/2}=f(2^k)\le M\sum_{j=0}^{s}A^j. \tag{6}
\]

If \(A\le1\), the right side is at most \(M(s+1)\), which grows at most linearly in \(k\). The left side grows exponentially—a contradiction. Thus \(A>1\).

The geometric sum in (6) now gives
\[
2^{k/2}
\le \frac{MA}{A-1}A^s
\le \frac{MA}{A-1}A^{k\log_b2}.
\]
Taking \(k\)-th roots and letting \(k\to\infty\),
\[
\sqrt2\le A^{\log_b2}.
\]
Therefore \(A\ge\sqrt b\). Together with (5),
\[
\boxed{f(n)=\sqrt n\qquad(n\in\mathbb Z_{>0}).}
\]

\paragraph{4. Extend to rational numbers and verify.}

For positive integers \(m,n\), multiplicativity gives
\[
f(m/n)=\frac{f(m)}{f(n)}=\sqrt{m/n}.
\]
Evenness and \(f(0)=0\) complete the classification.

Conversely, \(f(x)=\sqrt{|x|}\) is multiplicative, has the required normalisations, and satisfies the corrected reciprocal condition with equality. Finally,
\[
\sqrt{|x+y|}
\le\sqrt{|x|+|y|}
\le\sqrt{|x|}+\sqrt{|y|},
\]
so it is subadditive.

\subsubsection{Final corrected statement}\label{final-corrected-statement}

Use the problem above with the reciprocal condition \textbf{\(f(x)f(1/x)\ge1\)}, not \(\le1\). With that correction, the classification and proof are complete.
\end{proof}

\paragraph{Takeaways.}
\begin{itemize}
\item **Check inequality directions carefully:** Reciprocal substitutions only create rigidity when the bounds oppose one another.
\item **Do not assume monotonicity:** Numeral expansions provide growth estimates without comparing \(f(a)\) and \(f(b)\) directly.
\item **Change the base strategically:** Binary expansions give one bound; base-\(b\) expansions give the reverse bound at \(b\).
\item **Amplify away constants:** Taking high powers eliminates fixed multiplicative losses.
\end{itemize}

\end{dlsummary}
